%%%% CAPÍTULO 2 - REFERENCIAL TEÓRICO

\chapter{Referencial teórico}\label{cap:referencialTeorico}

Este capítulo reúne apenas os conceitos usados para formular e analisar o protótipo. As decisões de implementação são apresentadas no \autoref{cap:proposta}, os procedimentos no \autoref{cap:materialemetodos} e os valores medidos no \autoref{cap:resultados}.

\section{Negociação algorítmica e HFT}\label{sec:hft}

Negociação algorítmica designa a produção automatizada de decisões ou ordens a partir de regras computacionais. HFT é uma classe desse campo associada a grande automação, horizontes curtos de decisão, elevado fluxo de mensagens e investimento em infraestrutura de baixa latência. Velocidade, por si só, não define um sistema completo: uma operação real também envolve dados de mercado, controle de posição e risco, roteamento de ordens e requisitos regulatórios \cite{Aldridge2013,HasbrouckSaar2013}.

O caminho crítico é a sequência entre a chegada de um evento relevante e a disponibilização de uma resposta. Em software de propósito geral, esse caminho pode atravessar controlador de rede, pilha de protocolos, cópias de memória, chamadas de sistema e escalonamento. Técnicas como \textit{kernel bypass}, ou desvio da pilha de rede do núcleo do sistema operacional, e \textit{busy polling}, ou consulta ativa de novos dados, reduzem algumas dessas etapas, mas não são usadas no protótipo deste trabalho \cite{Rizzo2012,Jeong2014}.

Latência é o intervalo necessário para completar uma operação. Vazão é a quantidade de operações concluídas por unidade de tempo. \textit{Jitter} é a variação da latência entre observações; uma forma simples de expressá-lo é a diferença entre o maior e o menor valor de uma amostra. Média e vazão não descrevem, sozinhas, o comportamento temporal: mediana, percentis e máximo ajudam a revelar a cauda da distribuição \cite{Jain1991}. Neste trabalho, ``latência interna'' designa apenas o intervalo medido dentro do circuito FPGA, enquanto ``latência fim a fim'' designa a ida e volta entre o ARM e o FPGA.

\section{Livro de ofertas e dados de mercado}\label{sec:livroOfertas}

Um livro de ofertas limitado, do inglês \textit{limit order book}, organiza intenções de compra e venda. A melhor compra, ou \textit{best bid}, é o maior preço comprador disponível; a melhor venda, ou \textit{best ask}, é o menor preço vendedor. A diferença entre esses preços é o \textit{spread}. A profundidade corresponde às ofertas disponíveis além do primeiro nível \cite{Gould2013}.

Um livro por ordem registra cada oferta e sua prioridade. Um livro por nível agrega quantidades que compartilham o mesmo preço. A segunda representação perde informações sobre ordens individuais e prioridade temporal, mas reduz o estado necessário. Em ambos os casos, eventos de inserção, alteração, cancelamento e execução modificam o estado do livro \cite{Gould2013,Cont2010}.

O desbalanceamento usado neste trabalho é a diferença entre a quantidade da melhor compra e a quantidade da melhor venda, uma simplificação de sinais baseados no estado do livro discutidos na literatura de microestrutura. Ele não representa, isoladamente, previsão de preço, intenção de negociação ou lucratividade \cite{Gould2013,Cont2010}.

\section{FPGA, SoC e HPS}\label{sec:fpgaSoc}

Uma FPGA é um circuito integrado reconfigurável composto por blocos lógicos, registradores, memórias e interconexões programáveis. Diferentemente de um processador que executa uma sequência de instruções sobre recursos compartilhados, a FPGA implementa caminhos de dados e controle específicos. Operações independentes podem ser realizadas em paralelo quando há recursos e temporização suficientes \cite{HarrisHarris2013,PattersonHennessy2017}.

Um sistema em chip, do inglês \textit{System-on-Chip} (SoC), reúne diferentes subsistemas em um dispositivo. O Cyclone V SoC da DE10-Nano combina lógica FPGA e um \textit{Hard Processor System} (HPS), isto é, um subsistema fixo com processadores ARM, controladores de memória e pontes para a lógica programável. A ponte HPS--FPGA transporta transações entre software e hardware; por isso, seu custo integra a latência do sistema, ainda que não integre a latência do núcleo lógico \cite{IntelCycloneV,TerasicDE10Nano,IntelSoCFPGADesignGuidelines}.

Determinismo, neste contexto, precisa ser delimitado. Um circuito síncrono pode executar um caminho com número fixo de ciclos para as entradas aceitas, desde que não haja espera variável. Isso não implica tempo constante para uma aplicação Linux que inclui acessos a barramento, consulta de filas e escalonamento \cite{HarrisHarris2013,PattersonHennessy2017}. O texto distingue essas duas escalas em vez de atribuir ao sistema completo uma propriedade observada apenas no circuito.

\section{Síntese de alto nível e MATLAB HDL Coder}\label{sec:hls}

HLS é a transformação de uma descrição algorítmica em uma implementação no nível de transferência entre registradores, do inglês \textit{Register Transfer Level} (RTL). O RTL explicita operações e transferências sincronizadas por relógio. Ferramentas de HLS podem gerar VHDL ou Verilog a partir de subconjuntos de C, C++ ou MATLAB, mas o projetista continua responsável por tipos, interfaces, memória, paralelismo e metas temporais \cite{MathWorksHDLCoder}.

O MATLAB HDL Coder aceita funções MATLAB e modelos Simulink compatíveis com síntese. Entradas e saídas precisam ter dimensões conhecidas; alocação dinâmica, cadeias de caracteres e diversas bibliotecas de software não possuem correspondência direta em hardware. Aritmética inteira ou de ponto fixo torna explícita a largura dos dados e evita a inclusão não planejada de operadores de ponto flutuante \cite{MathWorksHDLCoder,MathWorksFixedPoint}.

\section{FAST e mFAST}\label{sec:fast}

FAST, sigla de \textit{FIX Adapted for STreaming}, é um padrão de codificação para fluxos de dados financeiros. Ele usa modelos de mensagem e operadores como cópia e diferença para reduzir redundância entre campos sucessivos \cite{FIXFAST}. mFAST é uma biblioteca C++ que codifica e decodifica mensagens segundo esses modelos \cite{mFAST}.

O uso de FAST não torna um fluxo sintético equivalente a um fluxo de bolsa. Transporte, temporização, conjunto de eventos e distribuição estatística também compõem o comportamento dos dados de mercado \cite{Gould2013,Cont2010}. O protocolo de controle de transmissão (TCP) entrega um fluxo confiável e ordenado de bytes; multicast UDP, perdas e pacotes fora de ordem pertencem a outro regime experimental e não são reproduzidos pela comunicação local adotada \cite{RFC9293}.

\section{MMIO, Avalon-MM e filas circulares}\label{sec:mmio}

Entrada e saída mapeada em memória, do inglês \textit{Memory-Mapped Input/Output} (MMIO), expõe registradores de um periférico no espaço de endereçamento do processador. Uma leitura ou escrita do software torna-se uma transação no barramento. Avalon-MM é a interface de memória mapeada utilizada pelo Platform Designer da Intel para conectar componentes iniciadores e periféricos \cite{IntelAvalon,IntelPlatformDesigner}.

Uma fila circular armazena elementos em um vetor e usa ponteiros de início e fim com retorno à primeira posição \cite{Cormen2022}. As convenções específicas de fila vazia, fila cheia, posição reservada e ordem de publicação adotadas no protótipo são decisões de projeto e, por isso, são apresentadas no \autoref{cap:proposta}.

A consulta repetida de um estado é denominada \textit{polling}. Ela simplifica o controle, mas ocupa tempo do processador. Acesso direto à memória, do inglês \textit{Direct Memory Access} (DMA), permite transferir blocos sem uma operação da CPU para cada palavra, à custa de configuração e controle adicionais. MMIO e DMA não são alternativas universalmente superiores: a escolha depende do volume, do padrão de acesso e da latência buscada \cite{IntelSoCFPGADesignGuidelines,IntelSoCDMA2019}.
